% uid: ""
% school: "lycee"
% level: "3eme"
% section: ["Economie et Gestion"]
% chapter: "Statistiques"
% title: "Heures de révision et note : tableau à double entrée"
% difficulty: "Moyen"
% points: null
% tags: ["statistiques", "double-entree", "distributions-marginales", "point-moyen", "frequences"]
% provenance: "own"
% dressing: "tableau à double entrée"
% rung: 3
% batch: "2026-11-20-stat-3eco-A"

% ---STATEMENT---
On a interrogé $30$ élèves sur le nombre $X$ d'heures de révision avant un devoir et noté la note $Y$ obtenue sur $20$. Les effectifs $n_{ij}$ sont donnés dans le tableau :

\begin{center}
\begin{tabular}{|c|c|c|c|}
\hline
$X\backslash Y$ & $8$ & $12$ & $16$ \\ \hline
$2$ h & $4$ & $2$ & $0$ \\ \hline
$4$ h & $3$ & $7$ & $2$ \\ \hline
$6$ h & $1$ & $5$ & $6$ \\ \hline
\end{tabular}
\end{center}

\begin{enumerate}
  \item Compléter le tableau avec les marges (effectifs marginaux de $X$ et de $Y$).
  \item Calculer les moyennes $\overline{x}$ et $\overline{y}$ et donner les coordonnées du point moyen $G$.
  \item Donner la distribution des fréquences marginales de $Y$ (en fractions irréductibles).
  \item Quelle proportion des élèves ayant révisé au moins $4$ heures a obtenu au moins $12$ ?
\end{enumerate}

% ---SOLUTION---
\begin{enumerate}
  \item Marges : somme de chaque ligne (pour $X$) et de chaque colonne (pour $Y$).
  \begin{center}
  \begin{tabular}{|c|c|c|c|c|}
  \hline
  $X\backslash Y$ & $8$ & $12$ & $16$ & Total \\ \hline
  $2$ h & $4$ & $2$ & $0$ & $6$ \\ \hline
  $4$ h & $3$ & $7$ & $2$ & $12$ \\ \hline
  $6$ h & $1$ & $5$ & $6$ & $12$ \\ \hline
  Total & $8$ & $14$ & $8$ & $30$ \\ \hline
  \end{tabular}
  \end{center}
  Contrôle : $6+12+12=30$ et $8+14+8=30$ $\checkmark$
  \item \textbf{Rappel.}\par
  \begin{prop}[Moyennes marginales]
  $\displaystyle\overline{x}=\frac1N\sum n_{i\bullet}\,x_i$ et $\displaystyle\overline{y}=\frac1N\sum n_{\bullet j}\,y_j$, où $n_{i\bullet}$ et $n_{\bullet j}$ sont les effectifs marginaux.
  \end{prop}
  $$\overline{x}=\frac{6\times2+12\times4+12\times6}{30}=\frac{12+48+72}{30}=\frac{132}{30}=4{,}4$$
  $$\overline{y}=\frac{8\times8+14\times12+8\times16}{30}=\frac{64+168+128}{30}=\frac{360}{30}=12$$
  $$ \ib{$G(4{,}4\,;12)$}.$$
  \item Fréquences marginales de $Y$ : effectif marginal divisé par $N=30$.
  $$f(8)=\frac8{30}=\frac4{15}\qquad f(12)=\frac{14}{30}=\frac7{15}\qquad f(16)=\frac8{30}=\frac4{15}.$$
  Contrôle : $\dfrac4{15}+\dfrac7{15}+\dfrac4{15}=1$ $\checkmark$
  \item On se restreint aux lignes $X=4$ et $X=6$ : $12+12=24$ élèves. Parmi eux, ceux dont la note est $12$ ou $16$ : $(7+2)+(5+6)=20$.
  $$\frac{20}{24}=\frac56\approx\ib{$83\,\%$}.$$
  \textbf{Piège :} la proportion demandée est calculée sur les $24$ élèves qui ont révisé au moins $4$ heures, pas sur les $30$ élèves.
\end{enumerate}

\textbf{Idée :} les marges ramènent un tableau à double entrée à deux séries simples, dont on calcule les moyennes.
